\documentclass[11pt,a4paper]{article}
\usepackage[margin=24mm]{geometry}
\usepackage{amsmath,amssymb,mathtools,amsthm}
\usepackage{fontspec}
\DeclareMathSizes{10}{10}{7}{6}
\setmainfont{Times New Roman}
\setsansfont{Arial}
\setmonofont[Scale=0.82]{Consolas}
\usepackage{xeCJK}
\xeCJKsetup{CJKmath=true}
\setCJKmainfont[BoldFont=SimHei,ItalicFont=KaiTi,BoldItalicFont=SimHei]{SimSun}
\setCJKsansfont{Microsoft YaHei}
\setCJKmonofont[BoldFont=SimHei]{SimSun}
\renewcommand{\contentsname}{目录}
\renewcommand{\abstractname}{摘要}

\usepackage{booktabs,longtable,array,graphicx,microtype,enumitem,xcolor}
\definecolor{reportblue}{RGB}{28,65,97}
\usepackage[colorlinks=true,urlcolor=reportblue,linkcolor=reportblue,citecolor=reportblue,filecolor=reportblue]{hyperref}
\hypersetup{bookmarksdepth=3}
\usepackage{xurl}
\setlength{\emergencystretch}{6em}
\setlength{\parskip}{0.4em}
\setlength{\parindent}{0pt}
\renewcommand{\baselinestretch}{1.08}
\widowpenalty=10000
\clubpenalty=10000
\displaywidowpenalty=10000
\setcounter{tocdepth}{1}
\raggedbottom
\DeclareMathOperator{\Var}{Var}
\DeclareMathOperator{\E}{E}
\newcommand{\ind}{\mathbf{1}}

\usepackage{needspace}

\title{概率证明如何完成}
\author{Shuo Deng \and Kenneth W. Shum\thanks{通讯作者: \href{mailto:wkshum@cuhk.edu.cn}{wkshum@cuhk.edu.cn}.}}\date{第 28 版\\2026 年 9 月 15 日\\结果更新至 2026 年 9 月 14 日}
\hypersetup{pdfauthor={Shuo Deng; Kenneth W. Shum}}\begin{document}\maketitle
\emph{ProbabilityTheoryFormalization 项目评估}

\textbf{第 28 版 · 2026 年 9 月 15 日 · 结果更新至：2026 年 9 月 14 日}

\section*{\texorpdfstring{\protect\hypertarget{abstract}{}摘要}{摘要}}
人工智能代理怎样完成概率证明，并使结果能够供后续任务使用？本报告评价ProbabilityTheoryFormalization项目，结合以往审核与代码修订，以及十一个任务组三种工作配置的两批运行，考察数学成果、审核与协作的实际作用，以及完成判断的可靠性。\hyperlink{ref-1}{\mbox{[1,4-12,24,41]}}

实验比较规定流程A、可查阅资料并求助的Sol自主作者B，以及由Astra协调Sol作者和检查服务的配置C。按共同数学成果标准，第一批A、B各完成10/11组，C完成11/11组；第二批三者均完成11/11组。第二批原验收拒绝的三个交付，经代码与调用检查确认已完成要求的数学内容。A在两批中的求解用时中位数均高于B、C，未显示相应的任务覆盖优势。两批分别报告，并保留任务说明与运行安排的差异。\hyperlink{ref-13}{\mbox{[13,14,38,41,42]}}

数学案例解释了这些结果怎样形成：审核推动作者由原始矩推出中心矩界；协作者完成的反演证明进入唯一性及后续应用；旧Scheffé交付遗漏的尾部范围，在明确要求后的两次单独修复中得到补齐。历史代码则展示从概率模型推出密度、组合局部极限以及修订调用者的实际工作。历史标签增长与可比数学进展分别检查，保存代码的构建和语义检查提供另一组直接依据。\hyperlink{ref-4}{\mbox{[4-11,17-20,25,27,43]}}

第二批过程分析核对全部33项运行的4,054次实际工具调用，另列179次工具发现，并从每项运行的首次检索和首次诊断形成66条过程记录。分析连接资料、反馈、代码修订与最终使用，展示函数表示和索引的适配、局部试验的不同结局，以及已检查代码未能传回作者的断点。重复工具问题进一步形成启动信息与参数预检原型，并完成各自的功能验证。\hyperlink{ref-45}{\mbox{[45,46]}}

这些发现支持三项改进重点：审核覆盖公开定理承担的数学范围，协作检查成果是否传到使用者并进入证明，评价将等价实现的完成结果与实际投入一起报告。具体执行过程解释了工作怎样推进；同一任务的共同验收与历史对象核查，则为核对进展判断与所评价数学内容的对应关系提供依据。\hyperlink{ref-21}{\mbox{[21-23,40,42,43]}}

\textbf{阅读指引。}\hyperlink{section-1}{第1节}介绍项目与问题，\hyperlink{section-2}{第2节}说明评价方法，\hyperlink{section-3}{第3节}报告成果和投入。\hyperlink{section-4}{第4节}解释完整成果需要哪些数学工作，\hyperlink{section-5}{第5节}追踪检索与反馈怎样进入修订和实际使用，\hyperlink{section-6}{第6节}讨论评价依据并给出结论。详细案例与历史分析见\href{zh-appendix.pdf\#appendix-r}{附录R}，统一起点的方法与逐项记录见\href{zh-appendix.pdf\#appendix-s}{附录S}，展开案例的证据位置见\href{zh-appendix.pdf\#appendix-t}{附录T}。

\section[项目如何组织与维护证明]{\texorpdfstring{\protect\hypertarget{section-1}{}项目如何组织与维护证明}{项目如何组织与维护证明}}
\hypertarget{evaluation-questions}{}

本报告评价的对象是 ProbabilityTheoryFormalization：这个项目将教材数学转化为可以审核、供后续任务使用，并随依赖变化持续维护的形式化证明。构建成功、评审通过和教材任务完成彼此相关，却是不同的成果。报告围绕三个问题展开：

1. \textbf{哪些数学工作已经完成并可供使用？}对外提供的定理是否表达预期命题、完成要求的证明工作，并接通使用它所需的先前结果与调用代码？ 2. \textbf{审核、协助与流程控制怎样影响工作？}哪些可观察的决定促成证明修复、确认既有成果、留下范围缺口或中断交付？作出决定时，执行者掌握了什么信息？ 3. \textbf{关于完成和进展的判断有多可信？}哪些记录对应独立意见和可以比较的数学对象？据此得到的标签、代码检查和资源测量究竟说明什么？

以往审核与代码修订、A/B/C 任务对照分别回答这些问题的不同部分。\hyperlink{section-2}{第 2 节}说明两类材料的关系，不将它们合并成一个完成率。本节先解释项目文档规定的设计。实际数学工作、运行故障和维护效果的证据，来自后文具有明确时点的分析；不能仅因设计中存在某个组件，就认定它已经发挥了作用。\hyperlink{ref-1}{\mbox{[1,4-12,26,34,35]}}

\subsection[从概率教材到形式化证明库]{\texorpdfstring{\protect\hypertarget{section-1-1}{}从概率教材到形式化证明库}{从概率教材到形式化证明库}}
\href{https://github.com/Kind-NK-Hill/ProbabilityTheoryFormalization}{ProbabilityTheoryFormalization} 原名 ToyApollo，研究借助人工智能将概率教材形式化。它的目标是把教材定义、命题和论证转化为 Lean 声明与证明，既表达预期数学含义，也能供后续章节使用。Lean 检查形式化证明；Mathlib 提供可用的通用数学库。\hyperlink{ref-26}{\mbox{[26,34]}}

工作从教材材料开始。提供的原文按编号小节或习题部分组织，任务计划明确待形式化的命题和可使用的先前结果，包括教材中的显式引用。程序把原文、任务要求、允许的依赖和工作草稿汇集成任务包，交给证明作者。实现时，作者可以检索允许使用的库，并编写所需的辅助引理。\hyperlink{ref-34}{\mbox{[34,35]}}

项目把数学决策与软件检查分开。人工智能代理选择证明路线、编写代码并判断数学内容是否对应原文；Python 程序准备任务、调用检查并执行接纳代码的条件；Lean 检查代码；SQLite 记录任务状态以及保留的构建和审核证据。草稿与\textbf{正式证明库}分开保存，后者是项目的官方 Lean 证明文件。项目验收将已审核的工作正式纳入该库。\hyperlink{ref-34}{\mbox{[34]}}

本报告中的\textbf{代理}是 Codex 中的一次人工智能执行会话，具有指定模型、对话上下文和工具。\textbf{作者}编写证明，\textbf{助手}提供代码或数学分析，\textbf{评审者}检查命题与证明；这些都是人工智能角色，并非本报告的人类作者。

\subsection[编写、审核与接纳证明]{\texorpdfstring{\protect\hypertarget{section-1-2}{}编写、审核与接纳证明}{编写、审核与接纳证明}}
程序把教材原文、任务要求和允许使用的先前结果交给人工智能作者；作者编写 Lean 代码，并用 Lean 检查。随后，另一位人工智能评审者判断形式化命题与证明是否满足原文和任务要求。需要纠正的工作进入修复或诊断。在已审核代码进入正式证明库之前，程序还要确认通过意见对应即将接纳的代码与材料。这是项目文档规定的流程；后文案例说明实际执行在哪些地方发生了偏离。\hyperlink{ref-34}{\mbox{[34,35]}}

Lean 与人工智能评审者回答不同问题。Lean 检查证明是否符合写出的形式化命题；\textbf{语义评审}检查命题是否表达了要求的数学内容、关键论证是否确已证明而非被当作假设，以及列出的调用代码能否正确使用结果。仅通过 Lean 检查，不能证明代码与教材含义一致。

\begingroup\small\setlength{\tabcolsep}{3pt}\renewcommand{\arraystretch}{1.14}
\par\nopagebreak\medskip\noindent\begin{minipage}[t]{\linewidth}
\begin{tabular}{>{\raggedright\arraybackslash}p{0.2\linewidth}>{\raggedright\arraybackslash}p{0.3\linewidth}>{\raggedright\arraybackslash}p{0.44\linewidth}}
\toprule
\textbf{阶段} & \textbf{负责者} & \textbf{产生什么结果？}\\\midrule

{准备任务} & {程序} & {汇集原文、任务要求、允许的依赖和工作草稿。}\\
{编写与检查代码} & {人工智能作者，使用检索和 Lean 工具} & {形成候选证明（当前代码版本），或得到指导下一次修改的错误信息。}\\
{审核数学内容} & {由管理流程的代理或人工操作者调用的另一位人工智能评审者} & {给出关于命题、证明责任和列出调用者的意见。}\\
{修复或澄清问题} & {人工智能作者；必要时另设只读诊断者} & {修改证明代码，或先明确命题与证明路线。}\\
{接纳已审核代码} & {程序} & {确认有效意见对应当前代码和要求，将代码纳入正式库并记录操作。}\\
\bottomrule\end{tabular}\end{minipage}\par\medskip\endgroup
如果证明路线已获认可，剩余问题只是缺少引理等局部缺口，工作返回作者修复。如果出现原文不匹配、把关键结论藏进假设，或证明路线不清，则可能先需要\textbf{只读诊断}：由人工智能定位问题，再继续编写证明。诊断者不写证明，也不替代评审者。流程还禁止在语义评审失败后重提未修改的候选。第 1.4 节介绍按风险触发的\textbf{数学路线检查（Math Gate）}，它在继续编写之前检查拟定命题和证明路线。\hyperlink{ref-34}{\mbox{[34]}}

准备请求与执行请求也有区别。例如，\nolinkurl{review-now} 准备审核材料并指出下一步动作，管理流程的代理或人工操作者仍需实际调用评审者。类似地，助手写出的代码必须从其工作区导出，传给能够使用它的角色。服务调用结束，不一定意味着作者已收到可用结果，因此案例会核对实际返回和读取的内容。\hyperlink{ref-19}{\mbox{[19,34]}}

最后，程序检查意见是否对应当前候选、辅助代码、原文和审核要求；记录称之为\textbf{评审绑定}。它防止一个版本的意见被用于批准另一个版本。纳入代码的操作失败时，原正式文件会恢复；既有正式证明在重新审核中失败时，通常保留代码并记录修复需求。这些检查维护决定与证据的对应关系，但不能保证人工智能的数学判断正确，附录R.1.2的 L3 案例展示了这一限制。\hyperlink{ref-34}{\mbox{[34,39]}}

\subsection[完成证明意味着什么]{\texorpdfstring{\protect\hypertarget{section-1-3}{}完成证明意味着什么}{完成证明意味着什么}}
下文将任务要求完成的数学工作称为\textbf{原文义务}，将调用者使用定理时必须提供的假设称为\textbf{公开前提}。把必需的中间结论写成公开前提，就将这部分证明留给了调用者。\textbf{候选}指一个具体代码版本。下面用项目以往的一次修订说明这种区别。

历史反演任务 \nolinkurl{thm_9_5} 将区间概率与特征函数积分联系起来，说明有效的条件证明仍可能留下未完成工作。早期候选已有换元、逐点分情况、控制收敛与缩放。名为 \nolinkurl{h_spine} 的假设提供三项分析条件：可积性、交换积分的依据和 Dirichlet 积分极限。后续候选逐步把这些职责移入证明内部。\hyperlink{ref-1}{\mbox{[1]}}

\begingroup\small\setlength{\tabcolsep}{3pt}\renewcommand{\arraystretch}{1.14}
\par\nopagebreak\medskip\noindent\begin{minipage}[t]{\linewidth}
\begin{tabular}{>{\raggedright\arraybackslash}p{0.22\linewidth}>{\raggedright\arraybackslash}p{0.32\linewidth}>{\raggedright\arraybackslash}p{0.4\linewidth}}
\toprule
\textbf{修订阶段} & \textbf{已有或新增的数学工作} & \textbf{仍由调用者提供的工作}\\\midrule

{初始条件证明} & {在分析条件成立时完成主论证} & {可积性、积分交换和关键极限}\\
{下一候选} & {由界与有限测度证明乘积可积} & {积分交换和极限}\\
{再下一候选} & {证明非负截断参数下的积分交换} & {Dirichlet 积分极限}\\
{完整组装} & {由误差界证明极限，并组装已证成分} & {移除 \nolinkurl{h_spine}}\\
\bottomrule\end{tabular}\end{minipage}\par\medskip\endgroup
非负截断参数的限制很关键。最终证明让参数趋向正无穷，因此无需最初计划的、对全部参数成立的更强辅助命题。缩小内部引理的范围，避免了不必要工作，又不削弱教材目标。相反，把必需条件留在公开 \nolinkurl{h_spine} 前提中，就会使交付定理之外仍有实质工作。早期评审认可了中间进展，而仅看一连串未通过标签会掩盖这些进展。\hyperlink{ref-1}{\mbox{[1]}}

\subsection[数学路线检查与任务调度]{\texorpdfstring{\protect\hypertarget{section-1-4}{}数学路线检查与任务调度}{数学路线检查与任务调度}}
\textbf{父任务}负责教材原始命题。\textbf{数学路线检查（Math Gate）}检查自然语言证明路线和拟定的 Lean 命题，用于原文不匹配、把核心证明工作转为公开假设，或反复失败却没有清晰路线的情形。多数小任务不使用它。检查给出 \nolinkurl{go} 时，可按计划继续父任务编写；给出 \nolinkurl{stop} 时，当前计划下的普通编写暂停，须先处理所指出的问题。暂停的具体范围取决于指令，不能仅凭这个标签推断所有支持工作都被禁止。\hyperlink{ref-26}{\mbox{[26,34]}}

对多个任务，批量规划读取记录状态与依赖，区分可开始编写、需评审或路线检查、以及被上游任务阻塞的工作。它可围绕下游依赖与潜在工作冲突组织工作队列。计划确定下一步，管理流程的代理调用相应角色，每个任务仍遵循各自验收路径。因此，普通作者队列为空，可能意味着下一步该由评审者或诊断者行动，并不一定表示全部完成或必须由用户介入。\hyperlink{ref-34}{\mbox{[34]}}

\subsection[库接口与验收后的维护]{\texorpdfstring{\protect\hypertarget{section-1-5}{}库接口与验收后的维护}{库接口与验收后的维护}}
同一教材概念在项目和 Mathlib 中可能有不同 Lean 表示。例如，导入以项目期望定义表述的定理，并不会自动把它转为关于 Mathlib 积分的定理。对重要共用概念，文档规定先引入教材定义及基本性质，再证明连接该接口与 Mathlib 的\textbf{转换定理}，后续工作通常使用 Mathlib 形式。这针对反复出现的数学接口，并非每个局部引理都必须经历的循环。既有可运行文件应予保留，除非具体复用问题要求修复接口。\hyperlink{ref-35}{\mbox{[35]}}

依赖选择以实际使用为依据。原文提到期望，只提供数学背景，不自动要求导入项目的期望定义。导入理由应是证明所需的声明、先前定理或转换结果，并记录于依赖链中。通用库基础设施是正当支持；但若某辅助定理承载了分配任务的数学内容，评审就必须将其作为该任务证明的一部分检查，不能因其名为“桥接”便豁免。\href{zh-appendix.pdf\#r1-1}{附录R.1.1}的方差展开比较与\href{zh-appendix.pdf\#r3-1}{附录R.3.1}的 Gamma/Dirichlet 历史通过具体证明分析这一区别。\hyperlink{ref-35}{\mbox{[35]}}

验收并非维护的终点。定义或定理签名变化，可能要求调用者提供不同参数或假设。任务和依赖记录识别受影响工作，构建与评审检查修改后的接口及其使用。SQLite 为底层证据保留可重建索引，并区分早期通过意见、对当前代码及支持的覆盖，以及仍需处理的工作候选。发生变化后，必须重新检查相关覆盖，而不只是保留通过标签。\href{zh-appendix.pdf\#r3-2}{附录R.3.2}的全变差修复具体展示了这一义务：原文的概率测度条件必须同时落实到定义与调用者。\hyperlink{ref-7}{\mbox{[7,34]}}

\section[评价证据与方法]{\texorpdfstring{\protect\hypertarget{section-2}{}评价证据与方法}{评价证据与方法}}
两类证据回答\hyperlink{evaluation-questions}{第 1 节}提出的项目问题。\textbf{以往审核及对应代码修订}展示项目发展中证明责任、对外命题、共享辅助代码和调用者怎样变化；还用于辨认独立判断、确认数学目标是否可比、检验失败后追踪中的选择效应，并直接检查保存代码。\href{zh-appendix.pdf\#r3}{附录R.3}及附录 \href{zh-appendix.pdf\#appendix-c}{C}、\href{zh-appendix.pdf\#appendix-d}{D}、\href{zh-appendix.pdf\#appendix-l}{L}分别给出这些分析自身的问题、方法和发现。\hyperlink{ref-1}{\mbox{[1,4-11,32]}}

\textbf{十一个任务组的 A/B/C 对照}，将项目规定的流程与两种替代组织方式放在相同任务上比较。交付和资源记录说明各次实际执行产出了什么、消耗了多少；执行记录解释具体修复、漏检和中断。\hyperlink{section-2-1}{第 2.1–2.4 节}说明实验设置，\hyperlink{section-3}{第 3–5 节}报告发现。实验不能替代关于项目维护的历史证据，历史中的通过也不计入实验成功数。\hyperlink{ref-12}{\mbox{[12-14,24,31,38]}}

\textbf{轨迹分析在两类材料中连接实际工作与结果。}它追踪待完成的数学责任、收到的意见或指令、可观察的回应，以及随后被评价的代码与交付。历史记录和实验日志的详细程度不同；无法据记录确认的接收、使用或动机，不作补推。\hyperlink{section-2-5}{第 2.5 节}说明重建方法与覆盖范围。本报告的总体评价问题、下文所述原实验动机，以及后来开展的回顾分析分别交代。\hyperlink{ref-9}{\mbox{[9,10,17-20,24,25,27,29,40]}}

\subsection[为什么比较 A、B、C？]{\texorpdfstring{\protect\hypertarget{section-2-1}{}为什么比较 A、B、C？}{为什么比较 A、B、C？}}
完成概率证明既要写代码，也要决定何时检查、怎样回应错误、何时提交。项目为这些决定规定了审核与修复步骤。实验考察这些安排能否帮助代理完整交付，并与两种自主组织方式比较：由 Sol 主作者自行安排，或由 Astra 协调 Sol 作者和检查服务。三者处理相同任务组。B、C 是实验中的替代方案，不是项目后续的生产阶段。比较指标包括完成情况、求解经过时间和记录到的模型活动。\hyperlink{ref-12}{\mbox{[12,14]}}

\textbf{A 按项目规定的流程工作。}Sol 作者选择证明路线、检索库和修改代码，其他 Sol 会话承担检查。项目流程规定审核与修复条件，以及代码进入正式库之前需要满足的要求。

\textbf{B 由 Sol 主作者自行安排工作。}它收到一般工作建议，可以查阅项目的流程和审核资料，也可以请求其他 Sol 会话协助编写或复核。主作者决定怎样使用这些资源、是否检查，以及何时提交；实验没有要求它走完 A 的规定流程。

\textbf{C 由 Astra 安排 Sol 的工作。}Astra 协调者可以读取任务和代码，给 Sol 作者具体的数学建议或修改指令，再根据作者的结果和检查意见决定下一步。它知道可用服务的用途和调用方法，但没有获得内部检查指令全文。它自行选择是否调用检查；检查者仍按后台预设要求工作。\hyperlink{ref-12}{\mbox{[12,18,38]}}

\begingroup\small\setlength{\tabcolsep}{3pt}\renewcommand{\arraystretch}{1.14}
\par\nopagebreak\medskip\noindent\begin{minipage}[t]{\linewidth}
\begin{tabular}{>{\raggedright\arraybackslash}p{0.17\linewidth}>{\raggedright\arraybackslash}p{0.24\linewidth}>{\raggedright\arraybackslash}p{0.3\linewidth}>{\raggedright\arraybackslash}p{0.23\linewidth}}
\toprule
\textbf{配置} & \textbf{谁安排工作？} & \textbf{可以获得哪些帮助？} & \textbf{检查怎样安排？}\\\midrule

{A：规定流程} & {Sol 作者在项目流程约束下工作} & {按角色提供项目指令和人工智能检查} & {流程规定审核与修复条件}\\
{B：Sol 自主组织} & {Sol 主作者} & {项目参考资料；其他 Sol 作者或复核会话} & {主作者决定是否调用及如何跟进}\\
{C：Astra 协调} & {Astra 协调者，Sol 作者实现证明} & {工具说明；Sol 作者和固定检查服务；返回意见} & {协调者决定是否调用及如何跟进}\\
\bottomrule\end{tabular}\end{minipage}\par\medskip\endgroup
这些角色都是人工智能代理。Sol、Astra 分别指记录中的 \nolinkurl{gpt-5.6-sol} 和 \nolinkurl{gpt-6-astra} 配置。证明作者和检查角色使用 Sol，只有 C 的协调者使用 Astra；均采用 medium 推理强度设置。medium 表示推理设置，不是模型能力等级。B 是能够查阅文档并请求协助的自主配置；实验没有另设只获得题目、既无项目资料也无其他代理帮助的最简基线。C 同时改变决策模型、角色结构和信息权限，因此比较对象是三套完整安排。\hyperlink{ref-12}{\mbox{[12,18,31]}}

在本次对照中，轨迹分析考察规定的检查或自选检查何时找出了缺失的数学工作、何时错误放行，协调者怎样使用返回反馈，以及工作为什么继续或停止。对同一任务，先核查各配置是否遇到相同困难，再追踪谁指出问题、意见是否送达、代码如何变化。附录R.1–R.2据这些过程解释实际安排的具体优点与失败，不把完成数量当作审核质量的直接度量。

\subsection[代理获得什么信息，怎样完成一次任务？]{\texorpdfstring{\protect\hypertarget{section-2-2}{}代理获得什么信息，怎样完成一次任务？}{代理获得什么信息，怎样完成一次任务？}}
每组获得教材原文、任务要求、允许使用的数学库和指定的先前结果。代理需要提交定理声明、证明、必要的辅助代码及说明，解释它们如何完成任务。本报告将这套代码和说明称为\textbf{交付}。\hyperlink{ref-12}{\mbox{[12]}}

一次运行由\textbf{外层实验执行代理}启动和管理，它也是 Codex 中的人工智能代理。A、B 或 C 在提供的材料和时间额度内组织求解，随后交付成果，由新的评审会话按共同标准评价。内部审核通过、代码进入项目正式库和实验评价通过是分别记录的事件，第 2.4 节说明评分方法。

外层实验执行代理与 C 的 Astra 协调者分工不同：前者管理实验运行，后者在 C 内部安排证明工作。外层代理可以干预任何一组，例如停止后续活动。第 2.3 节和具体案例分别说明这些操作及其影响。\hyperlink{ref-31}{\mbox{[31,38,39]}}

工作资料具有不同作用。\textbf{流程资料}说明项目怎样组织编写、审核和修复；\textbf{审核资料}说明检查时关注哪些问题，例如是否增加了题目没有要求的假设、是否完整证明了所需结论。\textbf{工具说明}解释可以请求什么服务以及怎样调用；\textbf{检查反馈}则指出本次代码存在什么问题，或给出通过意见。

A 的作者与检查者分别收到相应角色的项目指令。B 的初始提示给出规划、检索、检查和按需求助等一般建议，项目流程与审核资料供它主动查阅；文档可用不等于每次运行都读过。C 的协调者可以读取公开任务、代码和反馈，但未获得一般审核与数学检查的内部指令全文。反馈仍可能让它了解部分检查标准。因此，C 自主决定怎样组织工作；任务要求和后台检查要求并不由协调者任意制定。\hyperlink{ref-38}{\mbox{[38,39]}}

B 的主作者能够直接操作终端、编辑代码，也可以编写请求交给其他 Sol 会话。C 将协调和执行分开：Astra 通过工具读代码、指派作者并请求检查，Sol 作者操作终端和修改代码。协调者可以给出具体证明建议，但不能通过固定检查工具改写检查者的内部指令。两组的内部复核均属于求解过程，之后还要接受实验的共同评价。\hyperlink{ref-38}{\mbox{[38,39]}}

研究选择十一个任务组，覆盖不同数学工作与组织问题。部分任务组要求相互连接的多个结果。L、M、H 是计划时的工作量分组，下面列出具体题目，后文用这些标识定位案例。\hyperlink{ref-12}{\mbox{[12,31]}}

\begingroup\small\setlength{\tabcolsep}{3pt}\renewcommand{\arraystretch}{1.14}
\begin{longtable}{>{\raggedright\arraybackslash}p{0.3\linewidth}>{\raggedright\arraybackslash}p{0.64\linewidth}}
\toprule
\textbf{组别} & \textbf{数学任务}\\\midrule\endfirsthead
\toprule
\textbf{组别} & \textbf{数学任务}\\\midrule\endhead
\midrule\multicolumn{2}{r}{\small 续下页}\\\endfoot\bottomrule\endlastfoot
{L1} & {构造新息，并证明其部分和形成鞅}\\
{L2} & {教材11.4方差可加性与11.5弱大数定律}\\
{L3} & {Scheffé 引理、密度收敛及相关分布结论}\\
{L4} & {全变差公式与分布计算}\\
{L5} & {构造最大耦合}\\
{L6} & {单传感器与双传感器线性最小均方误差}\\
{L7} & {向量按坐标收敛与连续映射}\\
{M1} & {停止过程与停止期望}\\
{M2} & {四阶矩条件下的强大数定律}\\
{M3} & {在共同概率空间上的 Skorokhod 表示}\\
{H1} & {特征函数反演、唯一性及两个应用}\\
\end{longtable}\endgroup
\subsection[哪些运行进入比较，执行条件有哪些变化？]{\texorpdfstring{\protect\hypertarget{section-2-3}{}哪些运行进入比较，执行条件有哪些变化？}{哪些运行进入比较，执行条件有哪些变化？}}
B、C 的早期试跑暴露过两类执行问题：已经承诺提供的资料或工具未实际提供；检查意见或助手代码已经生成，却未能正常返回给负责下一步决策的代理。执行程序随后得到修复。由于这些问题影响代理能够使用的信息，主要比较与早期试跑分开报告。\hyperlink{ref-31}{\mbox{[31,38]}}

本报告将从启动求解开始、包含作者修改、求助、检查和同一计时下继续工作的全过程计为\textbf{一次求解运行}。第一批开发记录中共有 41 次；原主表为每个任务组选取 A、B、C 各一次交付，共 33 项。A 包含九次较早运行和 H1、M3 两次较后运行；B、C 各采用十一次执行程序修复后的运行。早期 B、C 试跑用于分析故障及反馈使用。各次具体执行仍须检查，程序修复本身不能保证之后所有调用都成功。\hyperlink{ref-12}{\mbox{[12,13,31]}}

同一次运行可能留下多次交付，汇总时按记录顺序优先保留已有共同评价的完整交付，没有完整交付时保留已有共同评价的首份交付。由此，A-H1 在原计时内继续完成的成果计入该次运行。该规则由事后汇总程序确认，尚未找到运行前固定它的记录。\href{zh-appendix.pdf\#n1}{附录 N.1}保留精确选择步骤与运行对应表。\hyperlink{ref-31}{\mbox{[31]}}

选入共同评价的十一个任务组中，三组的教材原文和任务单一致；允许使用的先前结果清单也一致，唯一例外是 A-L7 还列出一个没有实现的连续映射目标。求解期间的资料、工具和干预则有下列具体差别。\hyperlink{ref-31}{\mbox{[31]}}

\begingroup\small\setlength{\tabcolsep}{3pt}\renewcommand{\arraystretch}{1.14}
\par\nopagebreak\medskip\noindent\begin{minipage}[t]{\linewidth}
\begin{tabular}{>{\raggedright\arraybackslash}p{0.18\linewidth}>{\raggedright\arraybackslash}p{0.43\linewidth}>{\raggedright\arraybackslash}p{0.33\linewidth}}
\toprule
\textbf{执行差别} & \textbf{实际发生了什么？} & \textbf{如何解释结果？}\\\midrule

{B 的资料和工具} & {早期试跑未获得部分承诺的能力，后续补齐并验证可用} & {主比较采用补齐后的 B；早期试跑另列}\\
{B、C 的结果传递} & {后续程序修复了部分检查意见或助手代码无法返回的问题} & {分析以决策者实际收到的内容为依据}\\
{任务附加提示} & {B 的最终参考包移除了 L4 必须使用指定接口的约束，以及 M3/H1 模块名提示；A-L4 实际收到过相关条件指令} & {按具体提示报告差别；尚未量化这些提示的影响}\\
{A-L7 的外层中止} & {尚余约 61 分钟时，外层禁止继续编辑、求助和新审核} & {原交付反映提前中止；恢复后的成果单列}\\
\bottomrule\end{tabular}\end{minipage}\par\medskip\endgroup
C 扩展到不同任务时，也调整了允许启动的题目和时间额度设置；选入的 C 运行共用核心后台与角色提示内容。\href{zh-appendix.pdf\#o1}{附录 O.1}具体说明接口约束、组织线索、实际暴露情况，以及实现共同性的范围。这里保留影响解释的实际条件，后文结合代码和反馈判断其作用。\hyperlink{ref-38}{\mbox{[38]}}

9 月 13 日的 A-L7 续作在封存工作状态的副本上恢复原作者，并明确允许交付合理修正后的命题。续作交付计入第3.1节的完成结果，追加时间计入第3.2节的累计用时。原轮次记录与指标同时保留；此前41次运行的历史清单不因这次续作而改写。附录R.2.6说明续作条件与证明变化。\hyperlink{ref-39}{\mbox{[39]}}

第二批按预先登记顺序完成同样十一组的A、B、C各一次运行，共33项。教材、给定数学库、模型型号、推理设置和任务额度保持一致；五组任务追加了验收说明，A的部分复审与回修安排、部分参考资料以及客户端和工具调用方式也有调整。第一批已按修正标准重评，评分修正本身不是运行条件变化；提前给作者的说明则改变了求解输入。两批因此分别报告。具体差分见\href{zh-appendix.pdf\#q1}{附录Q.1}。\hyperlink{ref-41}{\mbox{[41]}}

\subsection[完成证据与提交后的共同评估]{\texorpdfstring{\protect\hypertarget{section-2-4}{}完成证据与提交后的共同评估}{完成证据与提交后的共同评估}}
判定实验任务完成，需要检查交付覆盖了哪些数学要求，不能仅凭命令成功或作者决定提交。Lean 检查实际写出的命题；语义评审判断它与任务要求的关系；项目在接纳代码前检查意见的对应关系和接纳资格。任务对照另由提交后的独立评审者，按下述共同程序作出报告使用的完成判断。因此，以下几类证据回答不同问题：

\begingroup\small\setlength{\tabcolsep}{3pt}\renewcommand{\arraystretch}{1.14}
\par\nopagebreak\medskip\noindent\begin{minipage}[t]{\linewidth}
\begin{tabular}{>{\raggedright\arraybackslash}p{0.3\linewidth}>{\raggedright\arraybackslash}p{0.64\linewidth}}
\toprule
\textbf{证据} & \textbf{回答的问题}\\\midrule

{构建} & {代码能否在记录的环境中通过检查？}\\
{语义评审} & {评审者是否认为命题与证明满足原文义务？}\\
{绑定与生产验收} & {项目是否针对所需代码及当前评审依据应用了合格意见？}\\
{提交后的共同评估} & {实验交付是否满足本次实验采用的标准？}\\
\bottomrule\end{tabular}\end{minipage}\par\medskip\endgroup
第3节采用两批一致的数学成果标准：交付覆盖完整目标，允许数学等价路线及可组合接口；可从原题设和允许库直接推出的冗余参数，通过独立调用检查确认。指定路线是否遵守另记。原共同评估提供已有证据，后续复核纠正不一致判定，原意见保留。以下介绍原评审程序；本版对66项的逐项复核规则与三项调整见\href{zh-appendix.pdf\#q2}{附录Q.2}。历史材料保持各自时点的含义，实验成果与纳入生产库仍分别记录。\hyperlink{ref-13}{\mbox{[13,26,42]}}

\textbf{共同评估}是本次实验在提交后进行的评价，有别于求解过程中的帮助与内部评审，也有别于生产验收。原始端点评估程序对原文、任务、候选、允许的上游材料及记录的技术检查，请求两份新的只读语义意见。评估者请求中不提供配置元数据。两票同为通过或失败即形成裁决；两票均无法判断则维持未决；票决不同最多触发一名新裁决者，读取相同材料和两份意见，解决实质分歧，而非按置信度选票。原始、一致性及定向端点调用均使用记录中的 Sol 中等推理配置。\hyperlink{ref-30}{\mbox{[30]}}；\href{zh-appendix.pdf\#n6}{附录 N.6}

任务组判定针对要求的完整交付。仅通过坐标收敛定理，不等于完成含两个目标的向量收敛组；对明确标为局部成果的评审，只评价该成果，不替代缺失的任务。结果格式错误也与数学问题分开处理。在表示定理案例中，评审者写出通过意见，却在必须填写数字的置信度位置填入对象；评估软件据此记录未决的执行结果，流程随后进入裁决。\hyperlink{ref-22}{\mbox{[22,28,30]}}

随后重新评估方差展开、Scheffé 引理及 Cauchy 范围中的解释不一致。每个候选在规定解读下分别获得初始意见。义务上的分歧、无效结果，或对各候选对应行为的不一致处理，都可能触发必要裁决。这些裁决者可以看到新的初始意见和跨候选一致性摘要，对这部分后续证据并不设盲。反演组在六份初审后，对称补充了原始提交说明，因为匿名、仅看代码的视角不足以判断文档要求。\hyperlink{ref-23}{\mbox{[23,30]}}

9 月 12 日比较的修正验收明确采用四项解释：方差可加性允许通用展开路线；Scheffé 引理采用任务中的尾部表述；Cauchy 任务以教材的标准 Cauchy 核心为准；连续映射任务允许对不足的假设作合理且已披露的修正。9 月 12 日 09:22，已有意见被复用，可确定的格式不匹配被修复，并按这些标准重新合成结果。连续映射的两份追加意见此前已在 05:17 的补充中取得。这是事后的标准修正，不是求解者新增证明。\href{zh-appendix.pdf\#r1}{附录R.1}说明数学依据并保留原判定。\hyperlink{ref-13}{\mbox{[13,15,22,30]}}

\href{zh-appendix.pdf\#n6}{附录 N.6}说明各阶段规则、评估输入和设盲限制。这里的\textbf{独立}指有记录的评审会话与角色分离；裁决者会有意接收先前证据，同一模型使用不同会话也不能证明模型错误相互独立。

9 月 13 日 A-L7 续作是另一次交付与评价。两份新评审分别判断已披露修正版命题的正确性，以及是否逐字满足原命题。附录R.2.6和\href{zh-appendix.pdf\#o2}{附录 O.2}报告其结果与追加投入；完成的续作计入第3.1节结果表，追加求解时间计入第3.2节累计用时。\hyperlink{ref-39}{\mbox{[39]}}

\subsection[如何追踪信息、决策与证明变化]{\texorpdfstring{\protect\hypertarget{section-2-5}{}如何追踪信息、决策与证明变化}{如何追踪信息、决策与证明变化}}
轨迹分析要解释任务怎样走到最终结果。对每个重点案例，我们连接数学要求、作者或协调者实际收到的信息、下一步行动、定理声明与证明的变化，以及最终接受评价的版本。历史案例也沿着这些问题追踪审核和代码修订，并额外检查前后记录中的数学目标是否可比。\hyperlink{ref-8}{\mbox{[8-11,17-20,25,27,29]}}

核心比较是关键事件前后的实际工作。审核可能促成新推导，也可能支持提交未修改的代码，或根本未能送达作者。协助者可能写出了有效证明，但成果没有进入最终交付。我们查看工具返回和阅读记录，对比代码，并追踪复用结果的实际调用，以区分这些过程。无法恢复的环节保持未确认；已有事件仍可解释其前后的工作。\href{zh-appendix.pdf\#r2}{附录R.2}呈现这些发现。\hyperlink{ref-17}{\mbox{[17-20,24,25,27,40]}}

案例来自已经调查的修复、协助与执行问题，支持对具体过程的解释。日志提取覆盖全部记录中的求解运行，但判断反馈是否正确、是否被采用，还需要进一步阅读。另一次专项分析对索引中的 127 项内部服务中的六项采用了同一套详细标准；其他案例各有记录明确的考察范围。这些材料尚不能给出全体正确反馈采用率或数学修复率。\href{zh-appendix.pdf\#n4}{附录 N.4}说明分析单位与覆盖，\href{zh-appendix.pdf\#g18}{G.18}呈现六项检查，\href{zh-appendix.pdf\#m2}{M.2}说明日志处理规则。\hyperlink{ref-16}{\mbox{[16,24]}}

另对A、B的两份旧L3代码，分别在原审核指令和明确尾部要求下各进行三次新评审，共12次；再对两份旧代码分别发起一次明确要求下的修复。诊断使用固定代码，修复使用另行冻结的指令，均与33项主运行分开。附录R.1.6报告结果，完整输入和判定见\href{zh-appendix.pdf\#q4}{附录Q.4}。\hyperlink{ref-43}{\mbox{[43]}}

第二批另按统一起点规则调查全部33项运行：每项取最早的数学检索及最早的证明或接口诊断，连接后续行动、局部成果与最终源码，共66条记录。25项运行的两个起点涉及同一需要，分别汇总。第5章报告这项过程分析，选择规则、逐项表及复核见\href{zh-appendix.pdf\#appendix-s}{附录S}。\hyperlink{ref-45}{\mbox{[45]}}

\Needspace{7\baselineskip}
\section[实验结果与投入]{\texorpdfstring{\protect\hypertarget{section-3}{}实验结果与投入}{实验结果与投入}}
两批运行已经结束。本章集中报告交付成果与投入；第一批工具记录补充说明求解中出现的问题。

\Needspace{7\baselineskip}
\subsection[两批任务的完成结果]{\texorpdfstring{\protect\hypertarget{section-3-1}{}两批任务的完成结果}{两批任务的完成结果}}
两批共66项主交付均已完成评价。按共同数学成果标准，第一批31项完成，第二批33项完成；下表逐项列出结果。

\begingroup\small\setlength{\tabcolsep}{3pt}\renewcommand{\arraystretch}{1.14}
\begin{longtable}{>{\raggedright\arraybackslash}p{0.22\linewidth}>{\raggedright\arraybackslash}p{0.113\linewidth}>{\raggedright\arraybackslash}p{0.113\linewidth}>{\raggedright\arraybackslash}p{0.113\linewidth}>{\raggedright\arraybackslash}p{0.113\linewidth}>{\raggedright\arraybackslash}p{0.113\linewidth}>{\raggedright\arraybackslash}p{0.113\linewidth}}
\toprule
\textbf{任务组} & \textbf{第一批 A} & \textbf{第一批 B} & \textbf{第一批 C} & \textbf{第二批 A} & \textbf{第二批 B} & \textbf{第二批 C}\\\midrule\endfirsthead
\toprule
\textbf{任务组} & \textbf{第一批 A} & \textbf{第一批 B} & \textbf{第一批 C} & \textbf{第二批 A} & \textbf{第二批 B} & \textbf{第二批 C}\\\midrule\endhead
\midrule\multicolumn{7}{r}{\small 续下页}\\\endfoot\bottomrule\endlastfoot
{L1} & {完成} & {完成} & {完成} & {完成} & {完成} & {完成}\\
{L2} & {完成} & {完成} & {完成} & {完成} & {完成} & {完成}\\
{L3} & {未完成} & {未完成} & {完成} & {完成} & {完成} & {完成}\\
{L4} & {完成} & {完成} & {完成} & {完成} & {完成} & {完成}\\
{L5} & {完成} & {完成} & {完成} & {完成} & {完成} & {完成}\\
{L6} & {完成} & {完成} & {完成} & {完成} & {完成} & {完成}\\
{L7} & {完成} & {完成} & {完成} & {完成} & {完成} & {完成}\\
{M1} & {完成} & {完成} & {完成} & {完成} & {完成} & {完成}\\
{M2} & {完成} & {完成} & {完成} & {完成} & {完成} & {完成}\\
{M3} & {完成} & {完成} & {完成} & {完成} & {完成} & {完成}\\
{H1} & {完成} & {完成} & {完成} & {完成} & {完成} & {完成}\\
{合计} & {10/11} & {10/11} & {11/11} & {11/11} & {11/11} & {11/11}\\
\end{longtable}\endgroup
第一批A、B的L3交付要求序列全部项可积，尚缺覆盖最终可积序列的尾部论证。

第二批原验收为30项通过，拒绝A-M1、B-M2和C-L3。统一复核逐项检查已有交付：A-M1的存在定理与停止期望结论可以组合，在原输入上覆盖三种要求的情形；B-M2已直接证明完整序列结论，区别在于没有遵守指定子列路线；C-L3的额外前提可以在同一输入上由通用库事实推出，独立调用检查通过。三项因此按共同数学成果标准计入完成，路线遵守另记。附录R.1.5说明每项依据，附录Q保留原判定和复核程序。\hyperlink{ref-41}{\mbox{[41,42]}}

A的第一批L7计入后续获准完成的交付，第二批L3计入工程恢复后的交付；两者额外求解窗口均计入下一节投入。另有两份旧L3候选的定向修复，均已通过，作为修复能力的独立证据报告于附录R.1.6。

\Needspace{7\baselineskip}
\subsection[同题比较求解用时]{\texorpdfstring{\protect\hypertarget{section-3-2}{}同题比较求解用时}{同题比较求解用时}}
求解时间从首次角色派发计至最后角色返回，包含运行内等待；并行协助的重叠时间只计一次。另行恢复的求解窗口累计计入，窗口之间的封存期及提交后的独立评价另列。两批各十一组的中位数如下。\hyperlink{ref-14}{\mbox{[14,41,42]}}

\begingroup\small\setlength{\tabcolsep}{3pt}\renewcommand{\arraystretch}{1.14}
\par\nopagebreak\medskip\noindent\begin{minipage}[t]{\linewidth}
\begin{tabular}{>{\raggedright\arraybackslash}p{0.26\linewidth}>{\raggedright\arraybackslash}p{0.22\linewidth}>{\raggedright\arraybackslash}p{0.22\linewidth}>{\raggedright\arraybackslash}p{0.24\linewidth}}
\toprule
\textbf{批次／累计分钟中位数} & \textbf{A} & \textbf{B} & \textbf{C}\\\midrule

{1} & {69.78} & {30.09} & {26.79}\\
{2} & {41.03} & {20.21} & {22.01}\\
\bottomrule\end{tabular}\end{minipage}\par\medskip\endgroup
逐项用时如下，单位均为分钟。每行对应同一任务组，可结合上一节的完成结果阅读。

\begingroup\small\setlength{\tabcolsep}{3pt}\renewcommand{\arraystretch}{1.14}
\begin{longtable}{>{\raggedright\arraybackslash}p{0.22\linewidth}>{\raggedright\arraybackslash}p{0.113\linewidth}>{\raggedright\arraybackslash}p{0.113\linewidth}>{\raggedright\arraybackslash}p{0.113\linewidth}>{\raggedright\arraybackslash}p{0.113\linewidth}>{\raggedright\arraybackslash}p{0.113\linewidth}>{\raggedright\arraybackslash}p{0.113\linewidth}}
\toprule
\textbf{任务组} & \textbf{第一批 A} & \textbf{第一批 B} & \textbf{第一批 C} & \textbf{第二批 A} & \textbf{第二批 B} & \textbf{第二批 C}\\\midrule\endfirsthead
\toprule
\textbf{任务组} & \textbf{第一批 A} & \textbf{第一批 B} & \textbf{第一批 C} & \textbf{第二批 A} & \textbf{第二批 B} & \textbf{第二批 C}\\\midrule\endhead
\midrule\multicolumn{7}{r}{\small 续下页}\\\endfoot\bottomrule\endlastfoot
{L1} & {85.82*} & {11.90} & {12.16} & {15.21} & {14.84} & {15.27}\\
{L2} & {14.25} & {30.09} & {18.00} & {16.75} & {16.30} & {18.98}\\
{L3} & {69.78} & {33.15} & {39.90} & {44.35} & {13.65} & {23.34}\\
{L4} & {46.11} & {15.89} & {16.54} & {41.03} & {18.79} & {18.04}\\
{L5} & {76.96} & {67.61} & {28.70} & {32.72} & {19.70} & {30.09}\\
{L6} & {34.34} & {21.98} & {23.61} & {24.51} & {34.13} & {21.49}\\
{L7} & {73.37} & {28.99} & {32.27} & {21.61} & {23.99} & {22.01}\\
{M1} & {136.89*} & {17.64} & {26.79} & {51.91} & {23.70} & {24.19}\\
{M2} & {59.54} & {36.38} & {36.53} & {51.17} & {40.56} & {22.91}\\
{M3} & {21.13} & {41.54} & {15.23} & {50.24} & {20.21} & {18.40}\\
{H1} & {182.16} & {126.17} & {103.36} & {135.59} & {123.89} & {86.42}\\
\end{longtable}\endgroup
*为由原始记录重建的估计。第一批A-L7为原运行58.86分钟加续作14.50分钟；第二批A-L3加计18.71分钟的工程恢复窗口。表中按未舍入秒数求和。完整端点与来源见附录M、P及Q.3。

同题比较比单个中位数更能说明差异。第二批B的中位数低于C，但B在5组更快，C在6组更快；H1分别为123.89与86.42分钟。A在两批均有较高的中位投入，但第一批L2用时最短。第一批L3还须结合未完成结果阅读。两批任务说明、部分流程和软件有变化，因此跨批耗时下降只作描述，不归结为某一个机制的提速。

\Needspace{7\baselineskip}
\subsection[工具问题在运行中的分布]{\texorpdfstring{\protect\hypertarget{section-3-3}{}工具问题在运行中的分布}{工具问题在运行中的分布}}
第一批全部33次选定运行都出现过命令不可用，27次出现文件访问问题。含Lean编译或接口报错的工具返回，A、B、C分别为180、291、322次。这些计数包括修改中的编译失败、接口试探和重复尝试。\href{zh-appendix.pdf\#m2}{附录 M.2}给出识别规则，\href{zh-appendix.pdf\#m3}{M.3}列出涉及的运行。\hyperlink{ref-24}{\mbox{[24]}}

一次任务可以经历多次报错和修改，最后仍然通过。A-M2就经历了多次构建失败，随后补齐证明并通过评价。因此，报错返回数衡量的是过程中的工具反馈，任务是否完成则由最终交付判断。附录R.2结合实际收到的反馈、代码修改和提交记录，分析哪些问题得到解决，哪些阻碍了交付。

\Needspace{7\baselineskip}
\subsection[模型用量与求解时间]{\texorpdfstring{\protect\hypertarget{section-3-4}{}模型用量与求解时间}{模型用量与求解时间}}
时间和词元用量反映不同投入。第一批M2中，B用时36.38分钟、记录输入15.237百万词元；A用时59.54分钟、记录输入13.880百万词元。两者均通过，用时较短的一方输入量反而较高。\hyperlink{ref-14}{\mbox{[14]}}

第一批词元比较使用三配置记录均完整的八组任务；A的L1、M1、L5仍有用量缺项。输入量已包含缓存输入，输出量已包含推理输出，均不重复相加。作者和内部协助、提交后的共同评估、后续重评与工程试探分别核算，详见\href{zh-appendix.pdf\#n5}{附录 N.5}。\hyperlink{ref-16}{\mbox{[16,31]}}

第二批记录输入合计：A为143.819百万词元，B为97.605百万，C为87.452百万；A-L3恢复另有6.071百万。B-H1存在未记录完整用量的容量错误调用，因此B总量为已知下界。资源数与时间并列解释，逐项用量与独立评价的已知词元小计见\href{zh-appendix.pdf\#q3}{附录Q.3}。\hyperlink{ref-41}{\mbox{[41]}}

\Needspace{7\baselineskip}
\section[数学工作怎样成为完整、可用的成果]{\texorpdfstring{\protect\hypertarget{section-4}{}数学工作怎样成为完整、可用的成果}{数学工作怎样成为完整、可用的成果}}
本章回答：完成标签背后，代理究竟完成了哪些数学工作？以下依次检查定理接收的输入范围、证明内部应补出的推导、协作成果的后续调用，以及历史修订中的数学发展。第一批运行、旧候选的单独修复和历史代码修订分别标明来源。第5章再从第二批的统一起点出发，追踪检索与反馈怎样进入具体修订。

\Needspace{7\baselineskip}
\subsection[范围遗漏怎样在审核中被接受，又怎样被补齐？]{\texorpdfstring{\protect\hypertarget{section-4-1}{}范围遗漏怎样在审核中被接受，又怎样被补齐？}{范围遗漏怎样在审核中被接受，又怎样被补齐？}}
第一批Scheffé引理任务（L3）用于后续概率密度问题。对非负函数序列，在给定的几乎处处收敛条件下，极限函数$f$可积，且函数积分趋于$f$的有限积分；目标是推出$\int|f_n-f|\to0$。争议发生在引理接收哪些序列：A、B原交付要求每一项都可积，共同数学成果标准则要求覆盖从某一项开始可积的序列。原任务说明使用“允许”尾部表述，后来共同评价明确采用这一范围；这一区别影响对当时作者和审核者决定的解释。\hyperlink{ref-23}{\mbox{[23,39]}}

两种输入范围确有差别。在$(0,1]$上令首项为$1/x$，以后各项及极限均为零。尾部每项可积，首项却不可积。因此，接受逐项可积输入的旧定理不能直接用于这条序列。补齐任务需要选择一个可积尾部，在尾部完成积分论证，再利用有限初始段不改变极限，把结论恢复到原序列。这里要增加的是一段处理允许输入的证明。\href{zh-appendix.pdf\#r1-2}{附录R.1.2}保留具体范围和代码依据。

\textbf{旧审核已经修复过其他问题。}A早期把极限函数的非负性与可测性作为额外前提。一般评审指出缺口后，作者在证明内部推出所需性质。但数学路线检查把逐项可积解释为有限实值积分的合法编码；一般复审随后接受了同一解释。它检查的下游概率密度本来就逐项可积，能够调用旧引理，所以没有检验更广的尾部输入。B自选的数学复核也看到了尾部版本，却把它作为可选推广，作者因此提交了原实现。\hyperlink{ref-29}{\mbox{[29,39,40]}}

这条过程说明了遗漏如何留下：局部前提已被修复，下游可以调用，审核据此接受了仍然较窄的引理。接纳程序确认意见对应正确代码和材料，并不再作一次数学范围判断；因而后续程序检查没有改变该决定。

为进一步检查这一判断，第二阶段固定两份旧代码，分别使用原指令和增加尾部范围核查要求的指令，各作三次评审。下面区分评审是否讨论缺口、是否要求修复，以及最后是否通过。\hyperlink{ref-43}{\mbox{[43]}}

\begingroup\small\setlength{\tabcolsep}{3pt}\renewcommand{\arraystretch}{1.14}
\par\nopagebreak\medskip\noindent\begin{minipage}[t]{\linewidth}
\begin{tabular}{>{\raggedright\arraybackslash}p{0.26\linewidth}>{\raggedright\arraybackslash}p{0.22\linewidth}>{\raggedright\arraybackslash}p{0.22\linewidth}>{\raggedright\arraybackslash}p{0.24\linewidth}}
\toprule
\textbf{固定来源与指令} & \textbf{讨论尾部遗漏} & \textbf{要求修复} & \textbf{判为通过}\\\midrule

{A，原指令} & {0/3} & {0/3} & {3/3}\\
{A，增加范围核查} & {3/3} & {3/3} & {0/3}\\
{B，原指令} & {3/3} & {1/3} & {2/3}\\
{B，增加范围核查} & {3/3} & {3/3} & {0/3}\\
\bottomrule\end{tabular}\end{minipage}\par\medskip\endgroup
A的差别首先在于是否发现较窄的输入范围。B的原指令评审都讨论了尾部表述，其中两份仍认为不必完成。增加明确要求后，两份代码的全部评审都要求修复。识别出可行的推广，与认定它属于当前交付义务，在这里是不同决定。这是固定案例、每条件三次的观察。

另行冻结指令的两次定向修复随后补齐了缺口。A加入存在截点的接口，B采用最终可积的表述并保留逐项可积的兼容定理；两者都只在可积尾部使用积分论证，再恢复原序列结论。它们分别用时7.28和6.98分钟，通过构建与共同验收。修复指令在诊断结果返回前已经固定，故这里展示的是明确要求下的新增实现。两次修复单列，第一批原交付仍按当时的代码计分。\href{zh-appendix.pdf\#r1-6}{附录R.1.6}及\href{zh-appendix.pdf\#q4}{Q.4}

\Needspace{7\baselineskip}
\subsection[反馈怎样使公开定理满足原来的数学条件？]{\texorpdfstring{\protect\hypertarget{section-4-2}{}反馈怎样使公开定理满足原来的数学条件？}{反馈怎样使公开定理满足原来的数学条件？}}
第一批四阶矩强大数定律任务A-M2提供了另一个前提变化的例子。教材假设独立变量具有共同均值$m$和统一原始四阶矩界$\mathbb E X_i^4\le c$。准备摘要却提到中心矩，初稿也要求调用者给出$X_i-m$的中心四阶矩界。要按教材条件交付，作者必须在内部证明从原始矩到中心矩的转换。\hyperlink{ref-25}{\mbox{[25]}}

作者先搜索强大数定律、独立性和积分乘积接口，阅读实现，并在小文件中核对声明类型。随后展开有限和的四阶矩，用独立性消去相应项，取得尾概率界，再应用Borel–Cantelli。最初八次构建中，前七次失败，第八次通过。此时中心矩路线已经实现，但公开声明仍把所需的中心矩界留给调用者。

一般评审指出这个差别。作者随后补出

\[
(x-m)^4\leq8(x^4+m^4),\qquad
\mathbb E(X_i-m)^4\leq8(c+m^4)
\quad\text{when}\quad\mathbb E X_i^4\leq c.
\]
第一条逐点不等式使原始矩界能够控制中心矩；第二条把这个控制带入已有证明。第十一版公开定理接收教材给出的原始矩条件，并调用此前完成的中心矩辅助结果。对应版本的复审通过，最终源码保留这一结构，提交后的共同评价也通过。\href{zh-appendix.pdf\#r2-2}{附录R.2.2}

\begingroup\small\setlength{\tabcolsep}{3pt}\renewcommand{\arraystretch}{1.14}
\par\nopagebreak\medskip\noindent\begin{minipage}[t]{\linewidth}
\begin{tabular}{>{\raggedright\arraybackslash}p{0.22\linewidth}>{\raggedright\arraybackslash}p{0.32\linewidth}>{\raggedright\arraybackslash}p{0.4\linewidth}}
\toprule
\textbf{阶段} & \textbf{调用者需要提供什么} & \textbf{作者已经完成什么}\\\midrule

{第八版构建通过} & {中心四阶矩界} & {基于中心矩展开的强大数证明}\\
{评审指出差别} & {教材只提供原始四阶矩界} & {明确尚缺的转换义务}\\
{第十一版交付} & {教材原始矩条件} & {内部推出中心矩界，再调用已有辅助证明}\\
\bottomrule\end{tabular}\end{minipage}\par\medskip\endgroup
这个前后对比解释了审核带来的实际变化。构建通过时，代码已证明自己写下的命题；补上转换后，公开命题才使用任务原有的条件。反馈之后发生了可定位的新推导、声明修改和复验，所以这里能够说明一项具体修复。教材、作者实现、工具返回和评审共同参与了这条过程，单看评审次数无法表达这项变化。

同类问题也会出现在时间索引上。第一批B-L1的初稿对所有自然数时刻要求可测、可积，而任务只给正时刻条件及$X_0=0$。收到建议后，作者在内部推出零时刻性质，修改对外接口，并保留原辅助引理。A-M2补的是一个界，B-L1补的是边界时刻的性质；两者都把应由证明完成的工作从调用者前提中移了出来。\hyperlink{ref-20}{\mbox{[20,27]}}

\Needspace{7\baselineskip}
\subsection[协作者完成的证明怎样进入后续任务？]{\texorpdfstring{\protect\hypertarget{section-4-3}{}协作者完成的证明怎样进入后续任务？}{协作者完成的证明怎样进入后续任务？}}
第一批反演与唯一性任务组H1包含五项相连目标：指数函数的一个界、特征函数反演、经反演证明唯一性、Cauchy样本均值应用和整数值判据。这里可以沿实际数学结论检查协助的作用。\href{zh-appendix.pdf\#r2-3}{附录R.2.3}记录A的同根完成过程及后续代码调用。\hyperlink{ref-17}{\mbox{[17]}}

作者最初要求助手处理整个反演证明，却把工作权限限定为只读。返回意见指出Dirichlet积分的缺口，没有可用补丁。作者随后请求积分极限与一致界，取得已构建的辅助代码并纳入草稿。另一位助手完成剩余反演证明后又遇到导出失败；外层执行环节恢复产物，作者才得以读取、应用、构建并请求内部评审。

因此，数学结果产生和作者能够使用它，发生在不同环节。取得的反演论证处理有限积分交换、振荡核的逐点极限、控制收敛以及端点的一半质量。设$\mu$为概率分布，$\phi_\mu$为其特征函数。对区间端点$a<b$，它得到

\[
\lim_{T\to\infty}\frac1{2\pi}\int_{-T}^{T}
\frac{e^{-iat}-e^{-ibt}}{it}\phi_\mu(t)\,dt
=\mu((a,b))+\tfrac12\mu(\{a\})+\tfrac12\mu(\{b\}).
\]
这一结果随后进入唯一性证明：比较特征函数相同的两个分布，在适当端点得到区间概率相等，再由测度确定论证推出同分布。Cauchy均值应用则计算独立和的特征函数，并调用新完成的唯一性定理确定分布。中间成果的贡献可以沿这两次实际调用定位。

\begingroup\small\setlength{\tabcolsep}{3pt}\renewcommand{\arraystretch}{1.14}
\par\nopagebreak\medskip\noindent\begin{minipage}[t]{\linewidth}
\begin{tabular}{>{\raggedright\arraybackslash}p{0.3\linewidth}>{\raggedright\arraybackslash}p{0.64\linewidth}}
\toprule
\textbf{已完成的成果} & \textbf{后续证明中的使用}\\\midrule

{反演定理} & {唯一性证明比较区间概率并确定分布}\\
{唯一性定理} & {Cauchy均值应用由特征函数确定分布}\\
{整数值判据所导入的唯一性文件} & {该分支实际采用指数实部的另一条论证，没有调用唯一性定理}\\
\bottomrule\end{tabular}\end{minipage}\par\medskip\endgroup
最后一行也有解释作用：文件依赖提供了可用环境，实际定理调用才说明某项数学结论被用了。这份代码中的整数值判据使用自己的证明路线，不能按导入次数给协助成果增加采用次数。

大范围协助同样有完成并被采用的实例。第一批B-M1将停止过程和停时期望整组交给支持作者，主作者取得补丁、整合、构建并请求数学评审。支持实现除了有限停止结果，还处理极限过渡。例如，对鞅$X$及可积停时$T$，若增量由常数$c$控制，

\[
|X_{T\wedge n}|\leq |X_0|+cT
\]
给出所需的可积控制。最终目标文件与已整合代码一致。第一批A-M3也采用了助手的整任务实现。与之相对，较早A-L5的助手虽产生已构建代码，返回的路径却在作者容器外，作者无法取得补丁而继续本地实现；现有证据没有确认该助手代码进入最终证明。\hyperlink{ref-18}{\mbox{[18,19,28]}}

这些例子分别展示只读建议、辅助引理和整组实现。协助的范围、权限与交付方式共同决定后续能做什么。第5.4节再用第二批的同类记录检查代码传回的具体断点。

执行恢复也可能用于交付已经存在的证明。第一批A-L7的最终定理签名与证明体在封存工作候选中已经存在；获准续作后，作者明确披露命题修正，完成审核和交付。因此，追加窗口完成的是已有证明的披露与验收。\href{zh-appendix.pdf\#r2-6}{附录R.2.6}保留续作前后的状态。\hyperlink{ref-39}{\mbox{[39]}}

\Needspace{7\baselineskip}
\subsection[历史修订怎样增加能够使用的数学内容？]{\texorpdfstring{\protect\hypertarget{section-4-4}{}历史修订怎样增加能够使用的数学内容？}{历史修订怎样增加能够使用的数学内容？}}
以往代码修订补充了单次运行之外的观察。它们涉及不同版本的证明库和调用者，下面按实际增加的数学工作说明变化；这些历史记录单独分析，不增加第3章实验分母。

\textbf{从概率模型推出密度。}归一化Gamma变量到Dirichlet分布的早期实现使用公理桥接。后来公理被移除，却把归一化变量的密度直接定义成目标公式，原本应由模型推出公式的工作仍然缺失。后续修订建立乘积分布、归一化映射、单纯形上的支撑、投影密度和变量替换，才接通模型与密度公式。这个变化解释了为何单查公理名称不能衡量全部数学进展。\hyperlink{ref-4}{\mbox{[4]}}

所得结果又被两个后续任务使用，但这些任务仍各有坐标图或Beta分支等义务。再后来，已证密度结论被加入公开父定理，改变的是对外承诺的位置。模型到公式的推导、共享结果的提取和公开接口的补全，分别发生在历史中的不同时点。\href{zh-appendix.pdf\#r3-1}{附录R.3.1}

\textbf{把局部结果组成同一个全局对象。}矩母函数任务的路线需要复数带状区域上的解析控制、极限提取和特征函数识别。历史代码已有真实的局部结果，但不同紧致片段上的极限必须来自同一个全局函数，并且在交叠处相容。后续构造及相容性证明补上了整体组装；列出已有局部引理还不足以完成这一步。完整证明早于支持模块拆分，下游调用者后来虽只改动一个调用，上游依据已从私有公理变成内部证明。\hyperlink{ref-6}{\mbox{[6]}}；\href{zh-appendix.pdf\#r3-2}{附录R.3.2}

\textbf{使条件定理能够用于具体模型。}保存的集券例题研究从$N$种等概率券中收集$m_N=\lfloor(N+1)/2\rfloor$种，以独立几何阶段等待时间之和$T_N$表示所需时间。六月的条件性证明允许$N=2$，此时只收集一种券，唯一等待阶段的成功概率为1，整行总方差为零。该配置又通过方差恒等式要求归一化尺度的平方等于总方差，同时要求尺度严格为正。形式检查由此推出矛盾，说明该配置没有可行实例。八月另一份上游中心极限定理候选虽曾被拒绝，已经包含实质的Lindeberg与Lyapunov分支；剩余工作是将具体独立行连接到所用概率模型。完成这个连接后，集券应用仍须把实际均值、方差换成教材采用的渐近归一化。\hyperlink{ref-8}{\mbox{[8]}}

对方差为正的规模，令$Z_N$为使用真实均值和方差的标准化变量，教材表达与它的关系为

\[
Z_N=\frac{T_N-\mathbb E T_N}{\sqrt{\operatorname{Var}(T_N)}},\qquad
\frac{T_N-N\log2}{\sqrt{N(1-\log2)}}=a_NZ_N+b_N.
\]
后续工作已经证明均值误差有界、方差具有所需极限，由此建立$a_N\to1$、$b_N\to0$，并在变化的仿射映射下转移分布收敛。它们与上游极限定理、具体分布连接各是不同成果。这里分析的是保存的独立阶段等待模型，材料没有另行交付它与原始逐次抽券停止时间的等价证明。\href{zh-appendix.pdf\#r3-3}{附录R.3.3}

任务前提改变还会传到定义和调用者。全变差修订加入原文的概率测度要求后，使用该表达式的位置必须提供相应证据。下游收敛谓词需要同时表达概率测度条件与距离收敛；若改写为条件成立时才要求收敛，就会在条件不满足时自动为真。实际修订传递所需证据，使调用者仍表达预期对象。\hyperlink{ref-7}{\mbox{[7]}}

这些历史解释了报告所说的可用成果：证明要来自所声明的模型，局部结果要组成目标结论，调用者也要接得上。第6.3节进一步说明，跨历史记录统计进展时怎样保持这些对象可比。

\Needspace{7\baselineskip}
\section[从检索、修订到实际使用]{\texorpdfstring{\protect\hypertarget{section-5}{}从检索、修订到实际使用}{从检索、修订到实际使用}}
第二批33项运行全部通过共同数学验收。本章进一步问：\textbf{检索与反馈怎样进入证明修订，并形成被实际使用的成果？} 相同的最终完成结果之下，代理经历了不同的资料查找、接口适配、局部试验和协作交付。以下先交代统一观察的范围，再沿同一证明需要连接资料、反馈、调整和结果。

\Needspace{7\baselineskip}
\subsection[沿同一证明需要追踪后续工作]{\texorpdfstring{\protect\hypertarget{section-5-1}{}沿同一证明需要追踪后续工作}{沿同一证明需要追踪后续工作}}
分析从11个任务组三种配置的完整工具记录开始。去重并配对后，共有4,054次实际工具调用，另有179次工具发现。记录保留调用输入和返回。一次调用可能同时搜索、修改和检查，所以还需读取内容，确认这些动作在处理哪一项证明需要。\hyperlink{ref-45}{\mbox{[45]}}

每项运行选择两个起点：首次主动查找数学名称、类型或源码，以及首次读入源文件后的证明或接口诊断。从起点向后，追踪同一需要下的查询、修改、检查结果和最终源码。若助手参与其中，还要核对作者收到的内容和采用位置。

\begingroup\small\setlength{\tabcolsep}{3pt}\renewcommand{\arraystretch}{1.14}
\par\nopagebreak\medskip\noindent\begin{minipage}[t]{\linewidth}
\begin{tabular}{>{\raggedright\arraybackslash}p{0.32\linewidth}>{\raggedright\arraybackslash}p{0.12\linewidth}>{\raggedright\arraybackslash}p{0.5\linewidth}}
\toprule
\textbf{观察材料} & \textbf{覆盖} & \textbf{本章用它回答什么}\\\midrule

{实际工具调用} & {4,054} & {哪些操作和返回实际发生，工具问题分布在哪里}\\
{首次数学检索及其后续过程} & {33} & {找到资料后怎样推进证明，与最终源码有何联系}\\
{首次证明或接口诊断及其后续过程} & {33} & {检查暴露了什么，后续如何处理同一需要}\\
\bottomrule\end{tabular}\end{minipage}\par\medskip\endgroup
两类起点共形成66条记录；25项运行的两个起点涉及同一需要，两类因此分别汇总。首次检索可以是正常准备，首次诊断也可能来自临时查询；本章从这些起点追踪资料怎样变得可用，以及出现的具体问题如何处理。这个规则覆盖每项运行的早期工作，后续新需要的全部过程尚未逐段标注。发现相关声明、局部检查通过和最终采用也分别记录：前者说明找到了待核实适用性的候选资料，后两者还需代码及其检查证据。选择规则与复核分工见\href{zh-appendix.pdf\#s1}{附录S.1–S.3}。本章展开的调用顺序与源码位置另列于\href{zh-appendix.pdf\#appendix-t}{附录T}。

\Needspace{7\baselineskip}
\subsection[找到相关定理以后，还需要哪些工作？]{\texorpdfstring{\protect\hypertarget{section-5-2}{}找到相关定理以后，还需要哪些工作？}{找到相关定理以后，还需要哪些工作？}}
33条检索记录与最终源码的联系如下。统计沿首次检索所处理的需要向后追踪，包括后续查找和修改；它描述这段工作的去向。

\begingroup\small\setlength{\tabcolsep}{3pt}\renewcommand{\arraystretch}{1.14}
\par\nopagebreak\medskip\noindent\begin{minipage}[t]{\linewidth}
\begin{tabular}{>{\raggedright\arraybackslash}p{0.79\linewidth}>{\raggedright\arraybackslash}p{0.15\linewidth}}
\toprule
\textbf{所跟踪过程与最终源码的联系} & \textbf{记录数}\\\midrule

{至少一项相关声明被显式使用} & {29}\\
{相关实现或部分成果保留} & {2}\\
{首位协作者的成果被后续交付替代} & {1}\\
{尚未确认联系} & {1}\\
{合计} & {33}\\
\bottomrule\end{tabular}\end{minipage}\par\medskip\endgroup
\textbf{分布函数资料怎样接入分位数证明。}第二批Skorokhod表示任务B-M3需要构造具有指定分布、且满足所需收敛性质的随机变量。支持实现从分布函数的下分位数入手，用分布函数的单调性、两端极限和右连续性建立分位数性质，再支持后续构造。

主作者首先查询分布函数和分位数名称，扩大搜索后委派支持实现。助手读取分布函数及相关表示的源码，找到了所需性质。问题出在两种函数表示：项目使用自己的分布函数表达式，库定理则以库中的分布函数为对象。初次失败的实现中，助手已经写出两者在每个实数处相等的引理。它随后尝试直接用\nolinkurl{change}把目标换成库表达式，检查仍报告两者并非定义上相同。

这里需要区分两个具体动作。逐点等式说明两个表达式取值相同；直接更换目标却要求系统把它们认作同一展开形式。助手后来利用已有逐点引理，通过函数外延性得到函数等式，再用该等式改写极限、单调性与右连续性目标。这样，库定理才能作用于原来的项目函数。

\begingroup\small\setlength{\tabcolsep}{3pt}\renewcommand{\arraystretch}{1.14}
\par\nopagebreak\medskip\noindent\begin{minipage}[t]{\linewidth}
\begin{tabular}{>{\raggedright\arraybackslash}p{0.21\linewidth}>{\raggedright\arraybackslash}p{0.73\linewidth}}
\toprule
\textbf{次序} & \textbf{记录中的行动或结果}\\\midrule

{查找并开始实现} & {主作者检索并委派；助手阅读相关源码，写出分位数基础及逐点相等引理}\\
{首次检查与直接转换} & {目标中的函数表示不匹配；直接更换表达式仍被拒绝}\\
{查找和修订} & {助手读取上确界接口、试查候选名称，并修改界的证明；函数表示错误仍在}\\
{明确改写并复验} & {从逐点相等得到函数等式，改写目标，同时完成其他界的修改；支持文件检查通过}\\
{核对最终实现} & {最终支持文件保留函数转换及分位数性质，供后续构造使用}\\
\bottomrule\end{tabular}\end{minipage}\par\medskip\endgroup
这个局部过程得到的是通过检查的分位数基础。助手返回时，完整Skorokhod构造仍待继续；整项任务的最终通过属于后续交付结果。初次逐点引理、失败的直接转换和成功的函数改写都保存在调用返回中，因此可以准确指出修订发生在哪里。上确界证明也在期间修改，整个文件通过构建不能全归于一个函数等式。

\textbf{审核者给出的路线怎样成为作者的证明。}第二批B-M1需要证明停止后的过程仍为鞅。审核者搜索后返回了停止下鞅的库接口，以及将下鞅与上鞅性质合成为鞅的接口。返回意见给出具体路线：先把原鞅视为下鞅并停止，再对负过程作同样处理，取负后得到所需上鞅性质，最后合并。作者实际收到这份意见，随后复读源码并写入对应实现。

第一次检查仍未通过。库接口允许含无穷值的停止时刻，当前任务使用有限索引；作者还需证明取最小值与转换回有限索引后的表达式一致。后续检索查找去除无穷值包装、强制转换和最小值的接口，试查中也出现不存在的候选名称。补齐显式索引等式后，检查通过。最终文件保留了停止下鞅、负过程和合成为鞅的调用。\href{zh-appendix.pdf\#t2}{附录T.2}列出意见送达、作者修订及最终使用的对应位置。

这两个过程中的资料都确实进入了证明。B-M3补的是函数表示之间的连接，B-M1补的是跨接口的索引关系。它们解释了检索结果何以变得可用；表中的29条显式使用则给出按统一起点规则观察到的覆盖。二者各有作用，当前记录没有据此计算各配置的整体检索效率。

\Needspace{7\baselineskip}
\subsection[诊断之后，原来的问题发生了什么变化？]{\texorpdfstring{\protect\hypertarget{section-5-3}{}诊断之后，原来的问题发生了什么变化？}{诊断之后，原来的问题发生了什么变化？}}
首次证明或接口诊断中，19项来自实际实现，14项来自临时声明查询或局部试验。临时查询可能同时返回可用声明的类型和其他名称不存在的信息；实现中的错误则可能要求修改待交付的证明。需要继续追踪相同目标，才能知道返回之后改变了什么。

\textbf{有限和及可测性问题逐步修复。}第二批A-L1构造新息及其部分和；证明需要让变量相对于观察历史可测，并保证部分和可积。首份实现同时遇到可测性改写、索引无法确定和有限和可积性名称不匹配等问题。作者先读取报错位置，查找有限和的积分与可积性声明，明确有限索引的类型并改用相应接口。随后还需把观察历史的投影表达式与当前函数连接起来，证明函数等式后才能使用已有的可测性结论。过程中出现的编辑错误和零项可积性目标也分别修复。最终同一任务文件通过检查，并保留上述可测性与有限和可积性证明。这里的进展由报错目标、对应修改和同一文件的后续检查共同确认。\href{zh-appendix.pdf\#t3}{附录T.3}

\textbf{一项局部试验没有走到相同终点。}第二批线性最小均方误差任务B-L6中，支持作者试图证明三个项的线性组合平方的积分恒等式。初次诊断显示乘积可积性调用未能匹配当前目标，平方展开还留下其他目标。助手尝试更换表达式，期间一次补丁命令不可用，之后错误再次出现。它继续查询常数乘法接口，尝试简化和直接转换，后续版本推进了乘积可积性，但没有出现整条三项平方积分恒等式通过检查的记录。

同一助手返回时，标明已检查的是高斯矩与独立性等部分。最终任务文件确有单传感器的二次表达式证明，整项运行也通过共同验收；这些证据各自对应自己的目标，不能作为那条三项试验已完成的依据。该试验在所跟踪记录中仍归为未见解决证据。

B-L6的首次检索过程则有可确认的采用。作者沿高斯分布及分布相同的表达继续查找，最终源码使用检索所得接口把已知分布的积分性质转移到噪声变量，证明均值为零和平方积分。返回资料中另一条方差接口没有被使用。因此，这项运行同时包含被采用的检索成果和没有确认完成的局部试验。\href{zh-appendix.pdf\#s2}{附录S.2}及\href{zh-appendix.pdf\#t3}{T.3}

A-L1的后续检查和同一文件支持原目标已修复；B-L6的最终任务完成，却没有补上那条局部试验的完成证据。追踪时同时核对目标身份和后续代码，才能保留这种差别。第3.3节的报错分布提供第一批背景，本节则解释第二批具体诊断之后的工作，完整逐项表见\href{zh-appendix.pdf\#s2}{附录S.2}。

\Needspace{7\baselineskip}
\subsection[局部检查通过的代码怎样到达作者？]{\texorpdfstring{\protect\hypertarget{section-5-4}{}局部检查通过的代码怎样到达作者？}{局部检查通过的代码怎样到达作者？}}
第二批A-M2和A-M3的首位支持作者都产生了通过局部检查的代码。主作者要继续工作，还需要从返回结果中取得这些代码，并在自己的工作区整合。下面列出这一传递过程的实际状态。

\begingroup\small\setlength{\tabcolsep}{3pt}\renewcommand{\arraystretch}{1.14}
\par\nopagebreak\medskip\noindent\begin{minipage}[t]{\linewidth}
\begin{tabular}{>{\raggedright\arraybackslash}p{0.33\linewidth}>{\raggedright\arraybackslash}p{0.61\linewidth}}
\toprule
\textbf{检查位置} & \textbf{A-M2与A-M3中观察到的状态}\\\midrule

{支持角色的工作区} & {已有通过局部检查的代码}\\
{首次返回} & {补丁为空}\\
{作者随后读取的文件} & {仍是任务输入}\\
{后续工作} & {重新委派后，由后续候选替换首支持候选，并以非空补丁传入作者工作区}\\
\bottomrule\end{tabular}\end{minipage}\par\medskip\endgroup
空补丁与作者文件共同确定了断点：首次局部成果没有通过这次返回进入作者文件。后续非空补丁来自另一候选，所以最终完成不能算成首份支持稿迟到后被采用。现有证据记录了传递结果，没有确定空补丁产生的内部原因。

两题的首支持候选均被后续候选替换。第5.2节检索表却只有一条被替换记录，因为A-M3的首次检索起于主作者更早的构造需要，其后续成果进入最终证明；这里被替换的是首支持候选。追踪的对象不同，数量也就不同。\href{zh-appendix.pdf\#s3}{附录S.3}保留补丁、文件对应关系与源码位置。

这一检查也说明协助成果的采用怎样判定：先核对实际返回，再看作者文件的变化，最后查找后续证明中的使用。第4.3节H1展示成果传入后接通数学依赖，本节则定位成果尚未传入的环节。

\Needspace{7\baselineskip}
\subsection[将工具问题转化为可检验的改进]{\texorpdfstring{\protect\hypertarget{section-5-5}{}将工具问题转化为可检验的改进}{将工具问题转化为可检验的改进}}
前几节沿证明需要读取过程。本节回到全部工具调用，检查重复出现的执行问题。共有173次调用出现命令查找失败，其中49次外层退出码仍为0。调用可能包含不止一个内部命令，实际返回因此比单一外层状态提供更多信息。

其中，对\nolinkurl{apply_patch}的调用有66次查找失败，分布在全部33项运行和64个会话。新的会话可能再次尝试同一不可用入口；这个分布使启动时提供准确信息成为值得检验的改进。命令查找失败也可能涉及名称或搜索路径，具体类别与来源见\href{zh-appendix.pdf\#s4}{附录S.4}。

据此完成了启动工具信息卡；另针对参数拒绝实现了执行前的参数检查。两个原型各有独立的功能验证。\hyperlink{ref-46}{\mbox{[46]}}

\begingroup\small\setlength{\tabcolsep}{3pt}\renewcommand{\arraystretch}{1.14}
\par\nopagebreak\medskip\noindent\begin{minipage}[t]{\linewidth}
\begin{tabular}{>{\raggedright\arraybackslash}p{0.2\linewidth}>{\raggedright\arraybackslash}p{0.28\linewidth}>{\raggedright\arraybackslash}p{0.46\linewidth}}
\toprule
\textbf{原型} & \textbf{要解决的问题} & \textbf{已完成的验证}\\\midrule

{启动工具信息卡} & {让作者在工作开始时看到实际可用的工具入口} & {在隔离容器中核对默认入口及搜索路径变化，确认信息卡能区分这些状态}\\
{参数预检} & {在请求执行前识别格式或取值不合规的参数} & {静态回放两个工具的3,896次输入，识别9次原有拒绝，接受3,887次合法参数，与原程序的参数条件一致}\\
\bottomrule\end{tabular}\end{minipage}\par\medskip\endgroup
信息卡验证确认了所提供工具信息的准确性，静态回放确认了参数条件与原程序的一致性。信息卡对代理尝试次数、总投入和数学成果的影响，留待固定条件的实际对照；参数是否合法则是执行前的独立检查。原型实现、功能验证及证据范围见\href{zh-appendix.pdf\#s4}{附录S.4}。

本章沿同一证明需要连接了检索、反馈、修订和使用，也定位了代码传递与工具执行的问题。相关资料能否接入当前目标、局部成果能否到达作者，解释了相同完成结果之下不同的工作过程。

\Needspace{7\baselineskip}
\section[讨论与结论：怎样组织数学工作]{\texorpdfstring{\protect\hypertarget{section-6}{}讨论与结论：怎样组织数学工作}{讨论与结论：怎样组织数学工作}}
前面的分析把交付结果、数学工作和执行过程连接起来。本章回到第1章的三个问题：怎样评价已经完成的成果，怎样确认审核与协作的作用，以及关于进展的判断是否建立在可比对象上。

\Needspace{7\baselineskip}
\subsection[完成结果与投入需要一起阅读]{\texorpdfstring{\protect\hypertarget{section-6-1}{}完成结果与投入需要一起阅读}{完成结果与投入需要一起阅读}}
共同数学标准使不同实现可以按同一任务目标比较。通过之后，证明内容仍有差别：第一批B在方差任务中提供了可选的自行展开路线，C在反演应用中覆盖了更广的Cauchy分布族。它们分别增加了推导和输入范围。检查这些内容，能够解释相同完成标签之下实际交付了什么。\href{zh-appendix.pdf\#r1-1}{附录R.1.1}及\href{zh-appendix.pdf\#r1-3}{R.1.3}；\hyperlink{ref-21}{\mbox{[21,23]}}

A在两批中的求解时间中位数均较高，未显示相应的任务覆盖优势。这支持继续检查规定流程消耗的投入。逐题结果也显示，中位数较低的配置并非每题都快，第二批B、C各有领先任务。这里的用时包含求解期间实际发生的协助、等待和修复；它衡量完成工作的实际投入。

两批的部分提示与运行安排有变化，配置也同时涉及模型和组织差别。对它们应用同一数学标准，使成果判定可比，但不会消除运行条件本身的差异。因此，用时先在每批内按同题比较，跨批观察用于说明结果是否重复出现；某个单独流程组件的提速效果仍需固定其他条件检验。\hyperlink{ref-41}{\mbox{[41,42]}}

\Needspace{7\baselineskip}
\subsection[审核和协作的作用可以落实到具体工作]{\texorpdfstring{\protect\hypertarget{section-6-2}{}审核和协作的作用可以落实到具体工作}{审核和协作的作用可以落实到具体工作}}
评价审核时，可以沿“指出什么—要求什么—改了什么”检查。第一批A-M2的意见指出原始矩与中心矩条件的差别，随后出现了对应推导和公开接口修改。旧L3则显示，注意到更广范围后，评审仍可能决定当前无需补齐。明确义务会影响验收决定；只统计意见数量或通过标签，无法区分这些作用。

协作评价可以沿“传回什么—怎样整合—用在哪里”检查。第一批H1的调用关系使支持证明对后续任务的贡献能够定位。第二批B-M1则将意见送达、作者适配和最终库调用连接起来；空补丁案例定位了已检查代码尚未进入作者文件的阶段。这些观察适合改进对应环节，但它们没有独自求解的反事实对照，不能换算为协作节省的分钟数。

任务范围和交付可用性因而是可直接落实的检查重点。补齐缺失工作、使成果到达使用者，都能形成可确认的贡献。审核还可以确认既有成果并支持提交：第一批主比较C-L1收到可读的通过意见后，提交了未修改候选。这是检查进入提交决定的实例。\href{zh-appendix.pdf\#r2-1}{附录R.2.1}

\Needspace{7\baselineskip}
\subsection[历史进展怎样与可比对象对应？]{\texorpdfstring{\protect\hypertarget{section-6-3}{}历史进展怎样与可比对象对应？}{历史进展怎样与可比对象对应？}}
第4.4节展示证明库在多次修订中增加了哪些数学内容。统计这种发展，还要确认前后意见评价的是不是可比目标，以及意见对应哪份代码。已有调查先选择意见直接对应预期候选的广义比较，再以有目的的核查选出数学目标也得到肯定可比判断的部分。同一任务可贡献多个前后比较关系，关系也可以共享端点；因此关系数与任务数分别报告，关系不能当作独立任务样本。\hyperlink{ref-10}{\mbox{[10,32]}}

\begingroup\small\setlength{\tabcolsep}{3pt}\renewcommand{\arraystretch}{1.14}
\par\nopagebreak\medskip\noindent\begin{minipage}[t]{\linewidth}
\begin{tabular}{>{\raggedright\arraybackslash}p{0.24\linewidth}>{\raggedright\arraybackslash}p{0.17\linewidth}>{\raggedright\arraybackslash}p{0.17\linewidth}>{\raggedright\arraybackslash}p{0.17\linewidth}>{\raggedright\arraybackslash}p{0.17\linewidth}}
\toprule
\textbf{历史材料范围} & \textbf{比较关系数} & \textbf{涉及任务数} & \textbf{前项通过数} & \textbf{后项通过数}\\\midrule

{广义直接绑定比较} & {461} & {248} & {116} & {324}\\
{目标肯定可比且对象资格满足的关系} & {128} & {79} & {28} & {45}\\
\bottomrule\end{tabular}\end{minipage}\par\medskip\endgroup
广义集合显示标签有明显增长。在目标也得到肯定可比判断的部分中，增长仍然存在。第二行是有目的核查的合格子集，并非第一行的随机样本；两行的趋势差不能全部解释成错误评价或目标变化。

\textbf{详细案例覆盖哪些历史变化？}详细异常调查还形成了另一份有目的选择的队列。它与上述461条关系重叠163条，通过数从32增至33；其余298条则从84增至291。详细案例擅长解释某种问题怎样发生，覆盖的增长却只是广泛记录的一部分。因此，这些案例不足以说明全部历史增长主要来自已发现的异常，也不能用来估计异常在所有任务中的发生率。\href{zh-appendix.pdf\#r3-5}{附录R.3.5}

\textbf{后续修订中的任务组成。}对重复修订的统计同样受任务组成影响。追踪先固定失败起点，再逐步核查后续目标和对应代码；“首步”相对于这个起点，一项任务可以贡献多个起点。在严格合格的材料中，首步通过为22/69，后续步骤为3/31。后续31步来自尚未通过且仍有合格证据的路径，涉及15项任务；这些任务对应的19个首步中，本来就只有1个通过。因此，进入后续追踪的任务在所观察首步中的通过比例已经较低。要检验追加修订的收益，需要保持任务与条件可比，不能把两个比例直接当成同一批任务的衰减曲线。\href{zh-appendix.pdf\#r3-6}{附录R.3.6}

\textbf{保存代码提供另一种直接检查。}已有研究固定14对文件，包括6对开发样本和另外选取的8对。28个文件中17个可编译，16个产出所需具名声明；这16个目标均通过保存的两条公理规则，另一个可编译文件缺少请求的声明。编译、具名目标存在和命题语义分别回答不同问题。例如，一份早期证明直接把结论当作前提返回；后期源码去掉该捷径，却在保存环境中编译失败。源码层面的改进与保存环境中的构建失败应同时报告。\hyperlink{ref-11}{\mbox{[11]}}

\textbf{额外规则比较的信息。}同一研究还比较了两份候选分别在两套规则下的结果。在已测的公理与直接前提属性上，规则相同或具有已知蕴含关系，额外组合没有比已利用这些关系的公平固定标准比较多回答问题。这是已完成的有界负结果。构造场景另用于验证程序，未增加真实任务成功。\href{zh-appendix.pdf\#r3-7}{附录R.3.7}保留抽样、全部状态和比较依据。

这也说明历史调查与第5章各自承担的任务。历史案例追踪较长的证明发展和调用者变化；第二批固定起点记录使每项运行的早期过程都得到覆盖。前者的材料选择与后者的统一起点不同，各自回答对应范围内的问题。报告将它们共同用于解释项目，但保持计数和结论的范围清楚。

\Needspace{7\baselineskip}
\subsection[结论与改进方向]{\texorpdfstring{\protect\hypertarget{section-6-4}{}结论与改进方向}{结论与改进方向}}
本报告表明，可供后续任务使用的概率证明，需要把任务原有条件、实际推导与调用者连接起来。A-M2将中心矩转换移入证明，H1把反演成果接入唯一性及应用，历史修订继续补齐模型、极限和调用接口。这些变化构成可以检查的数学进展。两批任务的完成结果与投入比较，则说明规定流程需要逐环节检验其收益。

审核与协作的作用落实在具体决定及其后续工作中：指出缺失范围、要求补出推导、确认已有成果、传回代码并完成整合。检索所得资料需核对适用性，并在必要时完成适配，才能成为当前证明的一部分。B-M3和B-M1的过程说明已有定理如何接通不同表示，A-L1记录了诊断后的实际修复，B-L6则保留了没有确认完成的局部试验。由此得到的过程解释，比工具调用次数更具体，也保留了不同尝试的真实结局。

据现有证据，改进可集中在三个环节：审核时逐项核对公开定理的范围；协助返回时检查补丁、接收文件及实际使用；启动会话时提供准确的工具入口信息。前两项对应已有的遗漏和传递断点，信息卡则已通过功能验证。评价这些改动时，继续使用共同数学标准，同时报告数学成果和包含等待、修复与协助的实际投入。信息卡的模型效果对照仍属后续工作。若进一步专门研究卡点后的检索与恢复，还需从已有日志中一致界定起点、失败、调整和终点，纳入成功、替换、放弃及无法确认的过程。

关于评价可信度，数学目标可比且对象资格满足的历史子集中仍有接纳增长；后续修订较低的通过比例伴随任务组成变化，不能直接解释为修订效果递减。保存代码检查进一步显示，构建与公理规则通过仍需配合对命题和推导的检查。这些已完成的分析支持全文的评价方式：同时说明证明了什么、成果怎样进入使用，以及哪些代码和记录支持进展判断。

\Needspace{7\baselineskip}
\section[参考资料与公开项目材料]{\texorpdfstring{\protect\hypertarget{section-7}{}参考资料与公开项目材料}{参考资料与公开项目材料}}
公开仓库提供流程实现、Lean 证明库、构建说明及选定历史案例，也包含一个隔离的\href{https://github.com/Kind-NK-Hill/ProbabilityTheoryFormalization/blob/d07f272850899b58612adf1c7dc202538503252f/docs/workflow_demo.md}{流程演示}，使用生产任务准备、构建、评审验证与应用函数。小型定义及调用者示例演示接口修复、调用者构建失败与过期评审拒绝。Lean 编译真实执行；默认语义意见是记录的教学评审，也可配置外部评审者。这是软件检查入口，与概率任务评估分开。\hyperlink{ref-36}{\mbox{[36]}}

下列数字引用指向技术附录相关章节。附录 C、D 保留历史方法与结果，G、L 保留详细轨迹和数学论证，M、N 保留测量与评价方法，O 解释配置差异、运行中断的原因和 A-L7 续作。\href{zh-appendix.pdf\#e2}{附录 E.2}说明哪些原始记录随完整资料包提供，哪些仍需访问档案。包内沿用的\href{EVIDENCE_GUIDE.pdf}{证据指南}和浏览器入口提供请求、意见、代码与来源映射；后续核验记录也随包保留。实验材料供本次私下评阅使用。

\begingroup\small\setlength{\tabcolsep}{3pt}\renewcommand{\arraystretch}{1.14}
\fontsize{9.4}{11.1}\selectfont\renewcommand{\arraystretch}{1.0}\begin{longtable}{>{\raggedright\arraybackslash}p{0.08\linewidth}>{\raggedright\arraybackslash}p{0.54\linewidth}>{\raggedright\arraybackslash}p{0.32\linewidth}}
\toprule
\textbf{引用} & \textbf{材料} & \textbf{技术附录}\\\midrule\endfirsthead
\toprule
\textbf{引用} & \textbf{材料} & \textbf{技术附录}\\\midrule\endhead
\midrule\multicolumn{3}{r}{\small 续下页}\\\endfoot\bottomrule\endlastfoot
{\hypertarget{ref-1}{}[1]} & {反演中的逐步数学进展} & {\href{zh-appendix.pdf\#l5}{L.5}；\href{zh-appendix.pdf\#a2}{A.2}}\\
{\hypertarget{ref-2}{}[2]} & {编辑、语义评审与注册} & {A、B}\\
{\hypertarget{ref-3}{}[3]} & {早期调度与代码编辑} & {A、\href{zh-appendix.pdf\#l1}{L.1}}\\
{\hypertarget{ref-4}{}[4]} & {Gamma 到 Dirichlet 的证明职责} & {\href{zh-appendix.pdf\#b4}{B.4}、\href{zh-appendix.pdf\#l2-1}{L.2.1}}\\
{\hypertarget{ref-5}{}[5]} & {证明义务与父任务重评} & {\href{zh-appendix.pdf\#b1}{B.1}-\href{zh-appendix.pdf\#b3}{B.3}、\href{zh-appendix.pdf\#l2-2}{L.2.2}}\\
{\hypertarget{ref-6}{}[6]} & {数学路线关卡与矩母函数证明组装} & {\href{zh-appendix.pdf\#b10}{B.10}、\href{zh-appendix.pdf\#l2-3}{L.2.3}}\\
{\hypertarget{ref-7}{}[7]} & {全变差及其使用者} & {\href{zh-appendix.pdf\#l2-4}{L.2.4}}\\
{\hypertarget{ref-8}{}[8]} & {中心极限定理与集券链} & {\href{zh-appendix.pdf\#l3}{L.3}}\\
{\hypertarget{ref-9}{}[9]} & {目标变化与评审对象} & {B、\href{zh-appendix.pdf\#l4}{L.4}}\\
{\hypertarget{ref-10}{}[10]} & {历史比较、失败起点与选择效应} & {C}\\
{\hypertarget{ref-11}{}[11]} & {保存代码检查、定义见证与规则比较} & {D}\\
{\hypertarget{ref-12}{}[12]} & {配置与共同任务指令} & {K、\href{zh-appendix.pdf\#n1}{N.1}}\\
{\hypertarget{ref-13}{}[13]} & {修正验收与当前比较表} & {\href{zh-appendix.pdf\#k8}{K.8}-\href{zh-appendix.pdf\#k9}{K.9}、M}\\
{\hypertarget{ref-14}{}[14]} & {逐次运行耗时、派发与词元} & {H、M、\href{zh-appendix.pdf\#n5}{N.5}}\\
{\hypertarget{ref-15}{}[15]} & {评估者修正与回归检查} & {\href{zh-appendix.pdf\#k9}{K.9}、\href{zh-appendix.pdf\#n6}{N.6}}\\
{\hypertarget{ref-16}{}[16]} & {反馈语义覆盖与证据索引} & {\href{zh-appendix.pdf\#g18}{G.18}、\href{zh-appendix.pdf\#n4}{N.4}}\\
{\hypertarget{ref-17}{}[17]} & {H1 实际复用、支持与恢复} & {\href{zh-appendix.pdf\#g13}{G.13}}\\
{\hypertarget{ref-18}{}[18]} & {B/C-M1 采用、检查与文档} & {\href{zh-appendix.pdf\#g11}{G.11}}\\
{\hypertarget{ref-19}{}[19]} & {协助产出、交付、恢复与采用边界} & {\href{zh-appendix.pdf\#g1}{G.1}-\href{zh-appendix.pdf\#g4}{G.4}、\href{zh-appendix.pdf\#g18}{G.18}}\\
{\hypertarget{ref-20}{}[20]} & {C-L1 可读评审与提交} & {\href{zh-appendix.pdf\#g8}{G.8}}\\
{\hypertarget{ref-21}{}[21]} & {L2 对应调用与一致解释} & {\href{zh-appendix.pdf\#g15}{G.15}}\\
{\hypertarget{ref-22}{}[22]} & {L7 控制、条件修正、证明与披露} & {\href{zh-appendix.pdf\#g16}{G.16}}\\
{\hypertarget{ref-23}{}[23]} & {L3 尾部与 H1 教材范围} & {\href{zh-appendix.pdf\#g13-1}{G.13.1}、\href{zh-appendix.pdf\#g17}{G.17}、\href{zh-appendix.pdf\#k8}{K.8}}\\
{\hypertarget{ref-24}{}[24]} & {原始工具返回、过程分类与共现} & {\href{zh-appendix.pdf\#m2}{M.2}-\href{zh-appendix.pdf\#m3}{M.3}}\\
{\hypertarget{ref-25}{}[25]} & {M2 原始矩到中心矩的证明修复} & {\href{zh-appendix.pdf\#g14-1}{G.14.1}；R02、R03、R14、R18}\\
{\hypertarget{ref-26}{}[26]} & {先前保存的本地架构、流程与生产约定} & {\href{zh-appendix.pdf\#a5}{A.5}、K、N；V01-V06、P01-P02}\\
{\hypertarget{ref-27}{}[27]} & {L1 同题数学路线、控制、信息与后续修订} & {\href{zh-appendix.pdf\#g5}{G.5}-\href{zh-appendix.pdf\#g10}{G.10}；N01-N13、R06-R08}\\
{\hypertarget{ref-28}{}[28]} & {M3 整任务采用与置信度格式裁决} & {\href{zh-appendix.pdf\#g12}{G.12}；U05-U06}\\
{\hypertarget{ref-29}{}[29]} & {L3 中间前提移除与绑定版本的修复} & {\href{zh-appendix.pdf\#g14-2}{G.14.2}；R02、R03、T02}\\
{\hypertarget{ref-30}{}[30]} & {共同评估规则、可见输入与实际会话设置} & {\href{zh-appendix.pdf\#n6}{N.6}；\href{zh-appendix.pdf\#e2}{E.2}；随包评估证据}\\
{\hypertarget{ref-31}{}[31]} & {运行选择、共用输入、实际额度与缺失用量} & {\href{zh-appendix.pdf\#n1}{N.1}、\href{zh-appendix.pdf\#n5}{N.5}；\href{zh-appendix.pdf\#e2}{E.2}；随包运行证据}\\
{\hypertarget{ref-32}{}[32]} & {历史目标决定与保存代码抽样规则} & {\href{zh-appendix.pdf\#c4}{C.4}、\href{zh-appendix.pdf\#d1}{D.1}；\href{zh-appendix.pdf\#e2}{E.2}；随包历史方法证据}\\
{\hypertarget{ref-33}{}[33]} & {新息评审恢复记录与未决归属} & {\href{zh-appendix.pdf\#g6}{G.6}；\href{zh-appendix.pdf\#e2}{E.2}；随包恢复证据}\\
{\hypertarget{ref-34}{}[34]} & {文档所述任务生命周期、修复、调度与正式状态维护} & {\href{zh-appendix.pdf\#a5}{A.5}、\href{zh-appendix.pdf\#e3}{E.3}；公开架构、流程、工作区状态与依赖文档}\\
{\hypertarget{ref-35}{}[35]} & {教材与 Mathlib 接口、依赖选择及证明责任} & {\href{zh-appendix.pdf\#e3}{E.3}；接口政策、依赖轨迹与语义评审标准}\\
{\hypertarget{ref-36}{}[36]} & {公开实现与隔离流程演示} & {\href{zh-appendix.pdf\#e3}{E.3}；架构与流程演示文档}\\
{\hypertarget{ref-37}{}[37]} & {候选状态与六次服务的来源对应} & {\href{zh-appendix.pdf\#n2-1}{N.2.1}、\href{zh-appendix.pdf\#g18}{G.18}、\href{zh-appendix.pdf\#e2}{E.2}；第 12 版证据指南}\\
{\hypertarget{ref-38}{}[38]} & {A/B/C 版本、工具与信息权限独立核查} & {\href{zh-appendix.pdf\#o1}{O.1}}\\
{\hypertarget{ref-39}{}[39]} & {9 月 13 日 A-L7 续作及 L3 审核责任核查} & {\href{zh-appendix.pdf\#o2}{O.2–O.3}}\\
{\hypertarget{ref-40}{}[40]} & {A-L3 审核与程序关口的逐步核查；C-L1 信息与提交序列} & {\href{zh-appendix.pdf\#o4}{O.4}; \href{zh-appendix.pdf\#g8}{G.8}–\href{zh-appendix.pdf\#g9}{G.9}}\\
{\hypertarget{ref-41}{}[41]} & {第二批运行、条件与投入} & {\href{zh-appendix.pdf\#q1}{Q.1}}\\
{\hypertarget{ref-42}{}[42]} & {共同成果复核与调用检查} & {\href{zh-appendix.pdf\#q2}{Q.2}}\\
{\hypertarget{ref-43}{}[43]} & {旧L3诊断及完成的修复} & {\href{zh-appendix.pdf\#q4}{Q.4}}\\
{\hypertarget{ref-44}{}[44]} & {第二批代码与反馈分析} & {\href{zh-appendix.pdf\#q5}{Q.5}}\\
{\hypertarget{ref-45}{}[45]} & {第二批工具反馈与证明过程研究} & {\href{zh-appendix.pdf\#s1}{S.1–S.3}；33项覆盖、固定双起点、来源与复核}\\
{\hypertarget{ref-46}{}[46]} & {启动信息与参数预检原型} & {\href{zh-appendix.pdf\#s4}{S.4}；命令查找失败、静态回放及后续对照方案}\\
\end{longtable}\endgroup
\end{document}\n